% =============================================================================
% Lecture 16: Trajectory Fine-Tuning
% Agentic Machine Learning Systems (Volume III), Chapter 16
% =============================================================================
% Thesis: training on trajectories teaches whatever the loss touches, so
% observations stay context, targets are chosen spans, and a lower loss is
% never by itself a release decision.
% Hook and synthesis: 2,100 recorded SWE-bench repair runs, in which the
% model wrote 12 to 31 percent of the transcript (MLSysIM
% Literature.SweBenchTrajectories; Trajectory Fine-Tuning, figure 2).
% Running case: the repair agent's recovery record trace_r19_01_redacted_v1
% (e2 refused, e3 schema query, e4 corrected call, e6 report), trained into
% trained_policy_v1 (Sec. 16.1 and 16.12).
%
% Build: make deck V=3 D=16_fine_tuning   (figures: images/make_figs.py)
% Plan:  lecture_map.yml in this folder.
% produced the real library output in images/data/.
% =============================================================================

\documentclass[aspectratio=169, 12pt]{beamer}
\usepackage{../../assets/beamerthememlsys}
\mlsyssetup{
  volume={Volume III},
  chapter={Chapter 16},
  chaptertitle={Trajectory Fine-Tuning},
  layer=learning,
  logo={../../assets/img/logo-mlsysbook.png},
  instlogo={../../assets/img/logo-harvard.png},
}
\usepackage[T1]{fontenc}
\usepackage[scaled=0.9]{helvet}
\usepackage{courier}
% Code and real tool output: Inconsolata is narrower than Courier.
\usepackage[varqu,scaled=0.95]{zi4}
\renewcommand{\familydefault}{\sfdefault}
\usepackage{booktabs,amsmath}
\usetikzlibrary{arrows.meta,positioning}
\graphicspath{{images/}}
\setcounter{mlsystotalsections}{5}

\title{Trajectory Fine-Tuning}
\author{Vijay Janapa Reddi}
\institute{Harvard University}
\date{}

% A large answer choice for predict and concept-question slides
\newcommand{\choice}[2]{\textbf{\textcolor{mlsysaccent}{#1}}\enspace #2}
% A choice that is highlighted from overlay #1 on (the answer)
\newcommand{\answerat}[3]{\alt<#1->{\colorbox{mlsysaccent!15}{\choice{#2}{#3}}}{\choice{#2}{#3}}}
% A goal or takeaway line with its thumbnail
\newcommand{\thumbline}[2]{%
  \begin{minipage}[c]{2.6cm}\centering\includegraphics[width=2.5cm,height=0.9cm,keepaspectratio]{#1}\end{minipage}%
  \hspace{0.3cm}\begin{minipage}[c]{0.74\textwidth}\small #2\end{minipage}}
% A takeaway line with a larger thumbnail
\newcommand{\thumblineL}[2]{%
  \begin{minipage}[c]{3.3cm}\centering\includegraphics[width=3.2cm,height=0.95cm,keepaspectratio]{#1}\end{minipage}%
  \hspace{0.3cm}\begin{minipage}[c]{0.70\textwidth}\small #2\end{minipage}}
% An identifier from the running case, in code type
\newcommand{\op}[1]{\texttt{#1}}
% An exact excerpt of real tool output, a config, or source code
\newenvironment{logblock}[1][7.2]{\par\ttfamily\fontsize{#1}{\dimexpr #1pt*6/5\relax}\selectfont\raggedright\setlength{\parskip}{0pt}}{\par}
\newcommand{\ind}[1]{\hspace*{#1\fontcharwd\font`0}}
% Golden-thread tag: one line on the goals slide naming the H·S·A term this lecture acts on
\newcommand{\hsathread}[1]{{\scriptsize\textbf{\textcolor{mlsysaccent}{H·S·A thread:}} #1\par}}

\begin{document}

% =============================================================================
% TITLE AND HOOK
% =============================================================================
\mlsysnote{[0.5 min] [CORE] Read the thesis aloud once. Lecture 15 decided which records may teach and which spans of each record are targets. Today turns one admitted record into a weight update and asks what the update actually teaches. Notes total about 70 minutes; the core lecture is about 53 minutes. FLEX: notes tagged [OPTIONAL] can be cut; cut first the per-call storage slide, the label-shift diagnosis, the gradient-accumulation incident, the packing configuration, the Super Tux Kart data, schema augmentation, the adapter record, the release scorecard, and the open questions. Every concept the chapter owns stays on a core slide.}
\mlsystitle{Trajectory Fine-Tuning}{Training on trajectories teaches whatever the loss touches, so observations stay context, targets are chosen spans, and a lower loss is never by itself a release decision.}{cover_sft_blueprint.png}

\begin{frame}{How much of a repair run did the model write?}
  \mlsysactivity{Predict}{2.5 min}
  \centering
  \mlsyssteps[width=0.94\textwidth,height=4.6cm,keepaspectratio]{roles_step}{2}\\[0.1cm]
  {\normalsize\only<1|handout:0>{Real repair agents fixing real GitHub issues. Guess the dark share.}%
  \uncover<2->{Model output: \textbf{12 to 31\%}. Only that region can hold a training target.}}
  \mlsyscredit{2,100 SWE-bench Verified trajectories (github.com/SWE-bench/experiments) via MLSysIM; Trajectory Fine-Tuning, figure 2. Characters, each message once.}
  \note{[2.5 min] [CORE] NARRATE: these are not hypothetical. Each bar pools the published transcripts of one real repair agent on SWE-bench Verified, a set of real GitHub issues in Python repositories: 500 runs for four configurations and 100 for one, 2,100 in all. Gray is the prompt (system message, demonstrations, task), green is everything tools and the environment returned, and the hatched region is what the model itself generated. ENGAGE: one minute; students write a guess for the hatched share. Most guess half or more, because they picture an agent as a model talking. Step 2: the model wrote 12 to 31 percent; observations are 56 to 72 percent. Characters stand in for tokens because the configurations use different tokenizers. WARN: students will read the dark region as the training target. It is only the region a target can come from; the lecture shows that the target is smaller still, and that what happens to the green region decides what the model learns. LINK: this is today's running problem; the synthesis returns to these bars. FLEX: [CORE].}
\end{frame}

\begin{frame}{Recall: what did curation hand to training?}
  \mlsysactivity{Retrieval}{1.5 min}
  \begin{columns}[c]
    \begin{column}{0.58\textwidth}
      {\Large\bfseries Question of the day}\\[0.2cm]
      {\normalsize The repair agent's recovery record: e2 refused, e3 schema query, e4 corrected edit, e6 report. Which responses did curation select?}\\[0.25cm]
      \uncover<2->{\small \textbf{\textcolor{mlsysaccent}{e3, e4, e6}} are targets. The refused e2 and every tool reply stay as context.}
    \end{column}
    \begin{column}{0.38\textwidth}
      \centering\safeimg[width=\textwidth,height=1.2cm,keepaspectratio]{thumb_recall.pdf}\\[0.3cm]
      \safeimg[width=0.42\textwidth,height=1.7cm,keepaspectratio]{hsa_locator_origin.pdf}\\[0.02cm]
      {\footnotesize Origin: training moves no exposure}
    \end{column}
  \end{columns}
  \mlsyscredit{Trajectory Fine-Tuning, \S16.1 (record \op{trace\_r19\_01\_redacted\_v1}, policy \op{recovery\_targets\_v1}). Locator: \S16.1.}
  \note{[1.5 min] [CORE] Retrieval from Lecture 15. Give 45 seconds. The record: the repair agent proposed a malformed edit (e2); the runtime validator refused it before any tool ran; the agent asked for the tool schema (e3), issued a corrected edit (e4), received a check summary (e5), and wrote a scoped final report (e6). Step 2: the target policy recovery\_targets\_v1 selects e3, e4, and e6. The thumbnail is that record with the three selected turns lit and e2 dashed. NARRATE the locator: only the origin dot is lit. Fine-tuning reshapes each model call but leaves every exposure (horizon, state, authority) to the runtime; no weight update widens or narrows what the agent may do. ENGAGE: why keep e2 in the record at all if it is never a target? (Expected: the correction e4 only makes sense after the refusal; e2 is the reason for e4.) FLEX: [CORE].}
\end{frame}

\begin{frame}{Why can a useful recovery record teach the wrong lesson?}
  \thumbline{thumb_homes.pdf}{Decide whether a behavior belongs in weights, context, or runtime.}\\[0.1cm]
  \thumbline{thumb_mask.pdf}{Rebuild each input as the model saw it, and place the loss on chosen spans only.}\\[0.1cm]
  \thumbline{thumb_pack.pdf}{Weight and pack variable-length examples without changing the lesson.}\\[0.1cm]
  \thumbline{thumb_bound.pdf}{Bound the cost of training only on recorded states, and say when labeling learner states pays.}\\[0.1cm]
  \thumbline{thumb_pref.pdf}{Choose the objective, behavior source, and parameter artifact, and price each.}\\[0.1cm]
  \thumbline{thumb_gate.pdf}{Release a fine-tuned policy on its closed-loop acceptance.}\\[0.05cm]
  \hsathread{moves no exposure; training changes the model}
  \note{[1 min] [CORE] The driving question is the title, word for word from the chapter. Six goals; the first four segments take one or two each, and the last segment takes the release. Each thumbnail is the shape of the figure that teaches it. Two terms for the whole lecture, glossed now: supervised fine-tuning (SFT) adjusts a model's weights to raise the probability of selected target tokens given the tokens before them; a policy is the model viewed as the thing that chooses the agent's next action. Running case: the repair agent's recovery record from the recall slide, trained into a candidate called trained\_policy\_v1.}
\end{frame}

% =============================================================================
% SECTION: RECORDS
% =============================================================================
\mlsysnote{[0.5 min] [CORE] First segment: from a stored record to the exact tokens a loss is computed on, and which of those tokens may move the weights.}
\mlsysdivider{Which tokens of a recovery record should move the weights?}{Trajectory Fine-Tuning, \S16.2 to \S16.4}
\section{Records}

\begin{frame}{A chat template fixes the delivered prefix, token for token}
  \begin{columns}[T]
    \begin{column}{0.58\textwidth}
      {\footnotesize\textbf{Real output: Qwen2.5's own template, one tool turn}}\par\vspace{0.05cm}
      \begin{logblock}[8.5]
        <|im\_start|>user\\
        Fix the connection-pool leak in src/pool.py<|im\_end|>\\
        \colorbox{BlueLine!12}{\textcolor{BlueLine}{<|im\_start|>assistant}}\\
        \textcolor{midgray}{[... read\_file call ...]}\\
        \textcolor{BlueLine}{<tool\_call>}\\
        \textcolor{BlueLine}{\{"name": "run\_tests",}\\
        \textcolor{BlueLine}{\ "arguments": \{"target": "tests/test\_pool.py"\}\}}\\
        \textcolor{BlueLine}{</tool\_call>}\colorbox{BlueLine!12}{\textcolor{BlueLine}{<|im\_end|>}}\\
        \colorbox{OrangeLine!15}{\textcolor{OrangeLine}{<|im\_start|>user}}\\
        \textcolor{OrangeLine}{<tool\_response>}\\
        \textcolor{OrangeLine}{FAILED test\_release: pool exhausted}\\
        \textcolor{OrangeLine}{</tool\_response>}\textcolor{midgray}{ [...]}
      \end{logblock}
    \end{column}
    \begin{column}{0.39\textwidth}
      \small
      \textcolor{BlueLine}{\textbf{Model output}} (blue): 51 of 304 tokens, from \colorbox{BlueLine!12}{assistant} to \colorbox{BlueLine!12}{its stop token}.\\[0.2cm]
      \textcolor{OrangeLine}{\textbf{Tool replies}} arrive in the \colorbox{OrangeLine!15}{\emph{user} role}.\\[0.2cm]
      \uncover<2->{Serialization records the template, tokenizer, and reply order with the example, or trains on a prompt the model never saw.}
    \end{column}
  \end{columns}
  \mlsyscredit{Our run: transformers 5.7.0 chat template of Qwen2.5-0.5B-Instruct, commit 7ae5576; system lines, the first call and the second reply omitted (images/tool\_runs.py). Trajectory Fine-Tuning, \S16.3.2.}
  \note{[2 min] [CORE] Real output: the chapter's two-call repair turn through a real model family's chat template, the function that turns messages into the token string the model trains and serves on; system prompt and tool schemas omitted. NARRATE: blue is what the model generated: two tool calls between tool\_call tags (the slide shows the second), ended by im\_end, its stop token. Orange is what the runtime returned. This template delivers tool replies inside a user turn; another family uses a separate tool role. Count: of 304 tokens, the model wrote 51, 16.8 percent; Step 2: the training example is this string, and the message list is only its source. The serialized dataset must record the tokenizer version, template, tool formatting, and response boundaries, because the same messages under another template are a different input. WARN: two errors students will make. Re-sorting tool replies into call order invents a prompt the model never saw when replies arrived out of order. And role tags carry no authority: if the model writes text that looks like a tool\_response tag, it is still model output. ENGAGE: if the runtime compacted history before this turn, what must the example contain? (Expected: the digest the model actually received.) FLEX: [CORE].}
\end{frame}

\begin{frame}{Weights hold procedure; context holds contracts; runtime enforces}
  \centering
  \mlsyssteps[width=0.94\textwidth,height=5.7cm,keepaspectratio]{homes_step}{3}
  \mlsyscredit{Trajectory Fine-Tuning, \S16.2.1 and table 1.}
  \note{[2.5 min] [CORE] Every response comes from two inputs: the weights and the context assembled for that call. Fine-tuning changes only the weights. Step 1: weights hold procedure, behavior stable across tasks: inspect diagnostics before editing, read the schema after a refusal. Step 2: context holds contracts, facts that vary across deployments, sessions, and turns: tool schemas, permissions, observations. A model that memorized an old schema keeps emitting the old argument name after the interface changes. Step 3: runtime holds enforcement: schema validation, authorization, sandbox containment, timeouts. A probability distribution over tokens cannot guarantee an invariant, so enforcement lives in neither the prompt nor the weights. LINK: this is the proposal boundary of Lecture 2: training changes what the model proposes; the runtime still decides what executes. Lecture 15's intervention ladder put weight changes after context, schema, and runtime fixes; this slide says what a weight change can hold once it is chosen. ENGAGE: a refund policy that changes weekly: weights or context? (Expected: context. Anything that changes faster than a train, evaluate, release cycle behaves as a contract.) WARN: ``fine-tuning moves the prompt into the weights'' is only true for the procedural part of a prompt. FLEX: [CORE].}
\end{frame}

\begin{frame}{Each call's delivered prefix is its own training example}
  \begin{columns}[c]
    \begin{column}{0.49\textwidth}
      \centering\safeimg[width=\textwidth,height=5.0cm,keepaspectratio]{triangle.pdf}
    \end{column}
    \begin{column}{0.49\textwidth}
      \centering\uncover<2->{\safeimg[width=\textwidth,height=5.0cm,keepaspectratio]{swe_tokens.pdf}}
    \end{column}
  \end{columns}
  \mlsyscredit{Left: Trajectory Fine-Tuning, \S16.3.3 (30 turns stipulated; derived). Right: 500 SWE-agent runs, Claude 3.5 Sonnet, SWE-bench experiments via MLSysIM; ratio derived.}
  \note{[2.5 min] [OPTIONAL] Left: a runtime that compacts history (Lecture 4) shows the model a digest of the earlier turns, so one concatenated transcript is no longer what any call saw. The loader must then store each call as its own example with its exact delivered prefix. If each turn adds about the same number of tokens, call $t$ carries $t$ turns, and $H$ calls store $H(H+1)/2$ turns against $H$ for one transcript: 15.5 times at 30 turns. This is the same triangle Lecture 2 priced at serving time, now paid in training storage. Step 2: real data. 500 SWE-agent runs record the input tokens sent on every call; summed over a run, the median is 387,196 tokens, against 3,852 output tokens, a median ratio of 102. Stored per call, model output is about 1 percent of all tokens (pooled over the 500 runs), and the targets are a part of that. ENGAGE: what serving-side change would let one transcript be a faithful example? (Expected: a runtime that only appends, never prunes or compacts; the pipeline must still verify it from the delivered-input manifests.) WARN: the token sums are as each agent recorded them; caching or truncation inside the agent is not visible here. FLEX: [OPTIONAL]; if cut, say the 15.5 times aloud on the next slide.}
\end{frame}

\begin{frame}{The loss mask trains one author and conditions on the other}
  \centering
  \mlsyssteps[width=0.96\textwidth,height=5.7cm,keepaspectratio]{mask_step}{3}
  \mlsyscredit{Trajectory Fine-Tuning, figure 1 and Definition 16.1, redrawn at slide size (simplified: the schema query e3 and its reply are omitted).}
  \note{[2.5 min] [CORE] Loss masking: the loss is computed only at selected positions; every other token stays in the input. Step 1: the serialized record in three lanes by author. Supplied: prompt, tools, task. Policy: the refused call e2 (dashed), the corrected call e4, the final report e6. Tools: the refusal and the tool result. Step 2: the mask $m_t$. It is 1 on e4 and e6 and 0 everywhere else, including e2, which the model wrote. Authorship does not make a target; curation does. Step 3: masked tokens are still read. Attention lets every later target condition on them, so gradients from e4 and e6 update the weights that encode the refusal. The model learns to act on a refusal without learning to write one. The masked loss, equation 1: the mean of $-\log \pi_\theta(x_t \mid x_{<t})$ over positions with $m_t = 1$. WARN: role confusion. Train on everything and the model learns that tool output continues its own stream; at serving time it keeps generating past its call and writes a fake tool result. The runtime discards those tokens, but the turn stalls and the output budget is spent. ENGAGE: e2 is masked; why not delete it from the input? (Expected: e4 would answer a problem the example no longer shows.) FLEX: [CORE].}
\end{frame}

\begin{frame}{Why does every target sit one token late?}
  \mlsysactivity{Diagnose}{2 min}
  \begin{columns}[T]
    \begin{column}{0.58\textwidth}
      {\footnotesize\textbf{transformers 5.7.0, loss/loss\_utils.py (excerpt)}}\par\vspace{0.05cm}
      \begin{logblock}[9.5]
        \ind{4}\# Shift so that tokens < n predict n\\
        \ind{4}labels = pad(labels, (0, 1),\\
        \ind{8}value=\textcolor{BlueLine}{\textbf{ignore\_index}})\\
        \ind{4}\textcolor{RedLine}{\textbf{shift\_labels = labels[..., 1:]}}
      \end{logblock}
      \vspace{0.3cm}
      {\small\renewcommand{\arraystretch}{1.3}%
      \begin{tabular}{@{}l*{4}{c}@{}}
        position & 1 & 2 & 3 & 4\\
        labels & \texttt{a} & \texttt{b} & \texttt{c} & \texttt{d}\\
        \textcolor{RedLine}{\textbf{shift\_labels}} & \textcolor{RedLine}{\texttt{b}} & \textcolor{RedLine}{\texttt{c}} & \textcolor{RedLine}{\texttt{d}} & \textcolor{BlueLine}{$-100$}
      \end{tabular}}
    \end{column}
    \begin{column}{0.39\textwidth}
      \small
      Position $t$ is scored on token $t+1$; $-100$ means no loss.\\[0.3cm]
      A loader that also shifts trains every target one token late.\\[0.3cm]
      \uncover<2->{\textbf{Double shift:} the library already shifts. Shift once, and check on a short example.}
    \end{column}
  \end{columns}
  \mlsyscredit{Source: github.com/huggingface/transformers, src/transformers/loss/loss\_utils.py, release 5.7.0 (Apache-2.0), ForCausalLMLoss, three lines excerpted; nn.functional shortened. Trajectory Fine-Tuning, \S16.4.2 and table 4.}
  \note{[2 min] [OPTIONAL] Real source from the most used training library, cut to the lines inside ForCausalLMLoss that matter; the loss is then fixed\_cross\_entropy(logits, shift\_labels, ..., ignore\_index). The small table shows the shift on four labels. Two lines carry the chapter's mechanics. Blue: ignore\_index $-100$ is the mask; a label of $-100$ contributes no loss, so masking is writing $-100$ into the labels of every non-target position. Red: the library shifts labels by one inside the loss, because the logits at position $t-1$ predict token $t$. ENGAGE: pairs, one minute: the loader also shifted. What happens? Step 2: every target and every mask is misaligned by one across the whole corpus. The loss still falls, so no dashboard shows it. The model learns to predict the token after each selected token, and the mask now covers the first token of each observation instead of the stop token. The fix is to shift in exactly one place and to decode a short synthetic example before launching. WARN: the end token belongs to the target. The authentic stop token closes each selected response with $m_t = 1$; without it, the model generates past its call until the output limit and the runtime marks the proposal truncated. FLEX: [OPTIONAL].}
\end{frame}

\begin{frame}{Observations fill 77.1\%; the loss sits on 4.9\%}
  \centering
  \safeimg[width=0.92\textwidth,height=3.2cm,keepaspectratio]{inventory.pdf}\\[0.25cm]
  {\normalsize Unmasked token average: each target's weight shrinks by\\[0.1cm]
  $\dfrac{\text{predicted positions}}{\text{targets}} = \dfrac{16{,}383}{800} = \mathbf{20.5\times}$\quad(7.4$\times$ with rationale selected)}
  \mlsyscredit{Trajectory Fine-Tuning, Napkin Math 16.1 (hypothetical inventory: illustrative); dilution derived in images/calcs.py.}
  \note{[1.5 min] [CORE] The chapter's inventory of one serialized example, illustrative: 16,384 stored tokens; 12,640 of them, 77.1 percent, are observations. Of 960 response tokens, 160 are the refused proposal and stay context, so 800 are targets, 4.9 percent of the input. If the reviewed rationale (1,400 tokens) is also selected, 2,200 targets, 13.4 percent. An unmasked token average divides by all 16,383 prediction terms (the first token has nothing before it), so each target's weight shrinks by 20.5 times, or 7.4 times with the rationale selected. The rest of the loss goes to predicting file paths, stack traces, and JSON. Doubling the observations raises the dilution to 36.3 times. WARN: masking saves almost no compute. Every masked token still passes through every layer, because later targets attend to it; only the output projection and the loss at masked positions can be skipped. The mask changes what is learned and leaves the cost unchanged. LINK: the hook's real runs had 56 to 72 percent observations in the transcript; stored one example per call, the share rises toward this inventory. FLEX: [CORE].}
\end{frame}

\begin{frame}{Vote: which spans of the record get $m_t = 1$?}
  \mlsysactivity{Vote}{2.5 min}
  \begin{columns}[T]
    \begin{column}{0.48\textwidth}
      \normalsize
      \choice{A}{every assistant turn: e2, e3, e4, e6}\\[0.2cm]
      \choice{B}{e3, e4, e6 and the schema reply to e3}\\[0.2cm]
      \answerat{2}{C}{e3, e4, e6, each with its stop token}\\[0.2cm]
      \choice{D}{only e4, the corrected call}
    \end{column}
    \begin{column}{0.48\textwidth}
      \small
      \uncover<2->{\textbf{A} trains the refused call: 160 of 960 response positions.\\[0.15cm]
      \textbf{B} trains the model to write tool results.\\[0.15cm]
      \textbf{D} drops two admitted lessons.}\\[0.3cm]
      \uncover<3->{\textbf{Cut off at the output limit?} Keep the valid prefix or mask the turn; never append a stop token.}
    \end{column}
  \end{columns}
  \mlsyscredit{Trajectory Fine-Tuning, Checkpoint 16.1, \S16.4.5, and Napkin Math 16.1 (illustrative); 16.7\% derived.}
  \note{[2.5 min] [CORE] ENGAGE: vote, then one minute with a neighbor, then vote again. Step 2: C. A is what a loader does when it selects every assistant message; in Napkin Math 16.1 that puts 160 of the 960 response positions, 16.7 percent of the response weight, on reproducing the malformed call, in every recovery record that has one. B selects an observation: the model learns to predict the schema instead of requesting it. D is safe but wastes the schema query and the report, both of which curation admitted. Each selected response keeps its own authentic stop token, so the model learns when to hand control back. Step 3: a response cut off by the output limit or a timeout. Two sound options: supervise only the valid prefix and append nothing, or mask the turn. A synthetic stop token teaches the model to end a partial command as if it were complete, which the runtime's truncation rule would otherwise catch. WARN: a library flag such as TRL's assistant\_only\_loss selects every assistant turn, which is answer A. FLEX: [CORE].}
\end{frame}

% =============================================================================
% SECTION: BATCHES
% =============================================================================
\mlsysnote{[0.5 min] [CORE] Second segment: a correct mask per example is not enough once examples of different lengths share a batch and a buffer.}
\mlsysdivider{How do long and short trajectories share a batch?}{Trajectory Fine-Tuning, \S16.5}
\section{Batches}

\begin{frame}{Gradient accumulation once disagreed with full-batch training}
  \begin{columns}[T]
    \begin{column}{0.53\textwidth}
      {\small\textbf{Hugging Face, ``Fixing Gradient Accumulation,'' 16 Oct 2024}}\par\vspace{0.1cm}
      \begin{mlsyscard}{midgray}\small
        ``... the correct loss should be computed by the \colorbox{GreenLine!15}{total loss across all batches} in a gradient accumulation step \colorbox{GreenLine!15}{divided by the total number} of all non padding tokens in those batches.''
      \end{mlsyscard}
    \end{column}
    \begin{column}{0.44\textwidth}
      {\small\textbf{Our run: two micro-batches}}\\[0.1cm]
      \small
      \begin{tabular}{@{}lr@{}}
        \toprule
        targets & 2,048 and 64 \\
        micro-batch losses & 0.541, 7.309 \\
        \midrule
        mean of the two means & \textcolor{RedLine}{\textbf{3.925}} \\
        global token average & \textcolor{GreenLine}{\textbf{0.746}} \\
        \bottomrule
      \end{tabular}\\[0.25cm]
      \small\uncover<2->{Sum the losses and the target counts, then divide once.}
    \end{column}
  \end{columns}
  \mlsyscredit{Quote: huggingface.co/blog/gradient\_accumulation. Run: transformers 5.7.0 ForCausalLMLoss, synthetic logits (images/tool\_runs.py). Trajectory Fine-Tuning, \S16.5.1.}
  \note{[2.5 min] [OPTIONAL] A real incident from October 2024. The post opens: ``Gradient accumulation is supposed to be mathematically equivalent to full batch training; however, losses did not match between training runs where the setting was toggled on and off.'' Gradient accumulation runs several micro-batches (slices processed one after another) before one optimizer step, and it should give the same update as one large batch. In the transformers Trainer it did not: each micro-batch's loss was a mean over its own targets, and the Trainer averaged those means. The fix, shipped within a day, sums the losses and divides by the total count of non-padding target tokens across all micro-batches. Right: our run of the library's loss on two micro-batches with very different target counts, 2,048 and 64, constructed so the first is well fit and the second is not. The mean of the two means is 3.925; the true token average is 0.746, 5.3 times smaller, because the 64 tokens of the second micro-batch got half the weight. Step 2: the rule for any reduction across micro-batches or data-parallel workers: sum numerators and denominators first. ENGAGE: when are the two answers equal? (Expected: when every micro-batch has the same target count; variable-length trajectories make that rare.) FLEX: [OPTIONAL]; the next slide carries the concept.}
\end{frame}

\begin{frame}{Token averaging gives the long trace 32$\times$ the lookup's weight}
  \begin{columns}[c]
    \begin{column}{0.50\textwidth}
      \centering\safeimg[width=\textwidth,height=5.0cm,keepaspectratio]{norm.pdf}
    \end{column}
    \begin{column}{0.47\textwidth}
      \small
      \textbf{Token averaging:} every selected token counts once.\\[0.15cm]
      \textbf{Example averaging:} every example's mean counts once.\\[0.35cm]
      \uncover<2->{Splitting one trace into $H$ per-call records gives it $H$ times the weight of an unsplit example under example averaging. Token averaging is unchanged.}
    \end{column}
  \end{columns}
  \mlsyscredit{Trajectory Fine-Tuning, equations 2 and 3, table 8, and Napkin Math 16.2 (illustrative); weights derived.}
  \note{[2 min] [CORE] Loss normalization: how examples of different lengths are weighted in the batch loss. A 30-turn repair trace shares a batch with a one-turn file lookup. Token averaging (blue) weights each example by its target count: 8,192 against 256 targets gives 96.97 against 3.03 percent, 32 times. Example averaging (orange) gives each example's mean loss the same weight, 50 percent each. Step 2: the choice interacts with serialization. Per-call serialization turns one trace into $H$ records. Under example averaging that trace's share rises from $1/B$ to $H/(B-1+H)$, $H$ times an unsplit example. Under token averaging nothing changes, as long as each turn's response is selected in exactly one record, because repeated prefixes repeat no targets. LINK: Lecture 15 specified the corpus mixture in selected target tokens; only token averaging carries that mixture into the update unchanged. ENGAGE: when would you pick example averaging? (Expected: when a few very long traces should not dominate the update.) WARN: neither is wrong. What is wrong is not declaring one, so that segmentation or micro-batching changes the mixture silently. FLEX: [CORE].}
\end{frame}

\begin{frame}{How much of the update is the repair trace?}
  \mlsysactivity{Estimate}{2.5 min}
  \begin{columns}[T]
    \begin{column}{0.40\textwidth}
      \begin{mlsyscard}{mlsysaccent}
        \small One repair trace: \textbf{6,000} targets.\\
        15 file lookups: \textbf{200} targets each.\\[0.1cm]
        Then split the trace into 20 per-call records of 300.
      \end{mlsyscard}
    \end{column}
    \begin{column}{0.56\textwidth}
      \small
      \renewcommand{\arraystretch}{1.25}
      \begin{tabular}{@{}lrr@{}}
        \toprule
        Trace's share & Token avg. & Example avg. \\
        \midrule
        one record & \uncover<2->{66.7\%} & \uncover<2->{6.25\%} \\
        20 records & \uncover<3->{66.7\%} & \uncover<3->{\textbf{57.1\%}} \\
        \bottomrule
      \end{tabular}\\[0.3cm]
      \uncover<3->{The data and the loss did not change; a serialization choice moved the trace from 6\% to 57\% of the update.}
    \end{column}
  \end{columns}
  \mlsyscredit{Trajectory Fine-Tuning, Exercise 3 (hypothetical inputs); shares derived in images/calcs.py.}
  \note{[2.5 min] [CORE] ENGAGE: pairs, 90 seconds, fill the four cells. Step 2: token averaging: 6,000 of 9,000 targets, 66.7 percent. Example averaging: one of 16 examples, 6.25 percent. Step 3: split the trace into 20 records of 300 targets, each turn selected once. Token averaging is unchanged, 66.7 percent. Example averaging now counts 20 of 35 examples, 57.1 percent: the same data and the same loss function, and the trace's influence rose about ninefold because of how it was serialized. WARN: students expect splitting to dilute the long trace; under example averaging it amplifies it. ENGAGE follow-up: the lookups exercise a capability the release must keep. Which pair would you choose? (Expected: token averaging if the corpus mixture was set in tokens; example averaging with one record per trace if each task should count once; either way, declare it and measure retention.) FLEX: [CORE].}
\end{frame}

\begin{frame}{Packing raises occupancy from 19.1\% to 76.6\%}
  \begin{columns}[T]
    \begin{column}{0.50\textwidth}
      \centering\safeimg[width=\textwidth,height=3.9cm,keepaspectratio]{occupancy.pdf}\\[0.15cm]
      {\small Eight examples, 512 to 32,768 tokens. Padding gives every row the longest length.}
    \end{column}
    \begin{column}{0.47\textwidth}
      {\footnotesize\textbf{TRL 1.16.0, SFTConfig (excerpt)}}\par\vspace{0.05cm}
      \begin{logblock}[6.9]
        packing: bool = field(default=False, ...\\
        \ind{2}"Whether to group multiple sequences into fixed-length blocks to improve computational efficiency and reduce padding."\\[0.1cm]
        padding\_free: bool = field(default=False, ...\\
        \ind{2}"... flattening all sequences in the batch into a single continuous sequence ... only supported with the FlashAttention 2 or 3 ..."
      \end{logblock}
    \end{column}
  \end{columns}
  \mlsyscredit{Left: Trajectory Fine-Tuning, Napkin Math 16.2 (illustrative); derived. Right: TRL 1.16.0, trl/trainer/sft\_config.py (Apache-2.0), help strings condensed.}
  \note{[2 min] [OPTIONAL] Sequence packing: several variable-length examples share one fixed-length buffer. Napkin Math 16.2: eight examples from 32,768 down to 512 tokens, 50,176 tokens in all. Padding every row to the longest allocates 262,144 positions, 19.1 percent occupied: the mean length divided by the maximum, so one long trajectory sets the waste. Packing puts the longest alone in one 32,768-token buffer and the other seven in a second: 76.6 percent occupied, 4 times fewer positions. Right: the real switch. TRL's SFTConfig exposes packing and padding\_free as one flag each, off by default; padding-free flattening requires a FlashAttention kernel that takes example boundaries. WARN: occupancy counts context and targets together. The speedup is a property of the attention kernel on the target accelerator and must be measured. ENGAGE: what can go wrong once two examples sit in one buffer? (Expected: the next slide.) FLEX: [OPTIONAL].}
\end{frame}

\begin{frame}{Packed examples must not read each other, and a test checks it}
  \begin{columns}[c]
    \begin{column}{0.60\textwidth}
      \centering\mlsyssteps[width=\textwidth,height=4.7cm,keepaspectratio]{blockmask_step}{2}
    \end{column}
    \begin{column}{0.37\textwidth}
      {\footnotesize\textbf{Our run: example 2's logits, packed vs alone}}\\[0.05cm]
      \footnotesize
      \begin{tabular}{@{}lr@{}}
        \toprule
        mask & max $|\Delta|$ \\
        \midrule
        ordinary causal & \textcolor{RedLine}{\textbf{1.20}} \\
        block-diagonal & \textcolor{GreenLine}{\textbf{0}} \\
        \bottomrule
      \end{tabular}\\[0.2cm]
      \small\uncover<2->{Perturb one example; every other example's logits must not move.}
    \end{column}
  \end{columns}
  \mlsyscredit{Trajectory Fine-Tuning, figure 3 and equation 4, redrawn. Run: random 4-layer Llama-architecture model, transformers 5.7.0 (images/tool\_runs.py).}
  \note{[2.5 min] [CORE] Packed-example isolation. Left, step 1: one buffer holding example 1 (5 positions) and example 2 (3), queries down, keys across. Under an ordinary causal mask, example 2's queries score example 1's keys: the red block is a leak. Step 2: the block-diagonal mask sets every cross-example score to minus infinity before the softmax; variable-length kernels take the boundaries as offsets. Labels need the same care: the last token of example 1 must not be trained to predict the first token of example 2. Right: our run on a small random model. With the ordinary mask, example 2's logits differ from its logits alone by up to 1.20; with the block-diagonal mask and positions reset per example, the difference is exactly 0. Keeping buffer-wide positions with this rotary-position model left a difference of $4\times10^{-7}$, floating-point noise, as the chapter's RoPE argument predicts. WARN: no loss curve shows a leak; the equivalence test is the only evidence. ENGAGE: why reset positions if RoPE is shift-invariant? (Expected: other position schemes are not shift-invariant, so only the test transfers across architectures.) Land the segment: packing is a memory layout; the mask, the label shift, positions, and the declared normalization must leave each example's lesson unchanged. FLEX: [CORE].}
\end{frame}

% =============================================================================
% SECTION: COVERAGE
% =============================================================================
\mlsysnote{[0.5 min] [CORE] Third segment: a perfect loader still trains only on recorded prefixes. The deployed policy acts on prefixes its own choices produce.}
\mlsysdivider{What happens at a state no record shows?}{Trajectory Fine-Tuning, \S16.6 and \S16.7}
\section{Coverage}

\begin{frame}{Exposure bias arises because teacher forcing hides mistakes}
  \begin{columns}[c]
    \begin{column}{0.55\textwidth}
      \centering\safeimg[width=\textwidth,height=4.6cm,keepaspectratio]{offcorpus.pdf}
    \end{column}
    \begin{column}{0.42\textwidth}
      \centering\safeimg[width=\textwidth,height=2.9cm,keepaspectratio]{navlab.jpg}\\[0.05cm]
      {\footnotesize Carnegie Mellon's Navlab 1 to 5. ALVINN (1988) learned to steer the first from recorded driving and met the same problem.}
    \end{column}
  \end{columns}
  \mlsyscredit{Diagram: Trajectory Fine-Tuning, figure 4(a) and Definition 16.3, redrawn. Photo: Firefly4342, CC BY-SA 4.0, Wikimedia Commons. ALVINN: Pomerleau 1988, via \S16.15.}
  \note{[2.5 min] [CORE] Exposure bias: training on reference prefixes and running on prefixes the policy's own outputs produce. Teacher forcing feeds the recorded history at every step, so the policy never sees the consequences of its own errors. Left: the corpus covers $s_1$ to $s_H$; at $s_2$ the learner deviates and walks into states no record shows. In the running case: trained\_policy\_v1 meets a version conflict instead of the malformed-argument refusal it was trained on, repeats its accepted edit without checking state, the runtime's version guard blocks it, and the trial exhausts its budget. Nothing in the corpus taught that correction. Right: the problem is older than language models. ALVINN learned road following from recorded human steering on Carnegie Mellon's Navlab vehicles (the blue van in the photo is Navlab 1); a car trained only on good driving never sees how to recover from drifting toward the shoulder. WARN: in text generation an unexpected token drifts the phrasing; in an agent an errant action changes the environment. ENGAGE: why do added recovery records help without settling it? (Expected: they cover only the recoveries someone recorded, and this policy reaches other states.) FLEX: [CORE].}
\end{frame}

\begin{frame}{What does a 2\% error cost over 30 steps?}
  \mlsysactivity{Predict}{2.5 min}
  \begin{columns}[c]
    \begin{column}{0.50\textwidth}
      \centering\mlsyssteps[width=\textwidth,height=4.9cm,keepaspectratio]{bounds_step}{2}
    \end{column}
    \begin{column}{0.47\textwidth}
      \small
      Error measured on \textbf{reference} states:\\ \hspace*{0.3cm}$\le \min\{H,\ H^2\epsilon\}$\\[0.15cm]
      Error measured on \textbf{learner} states, each slip costing at most $u$:\\ \hspace*{0.3cm}$\le \min\{H,\ uH\epsilon\}$\\[0.25cm]
      \uncover<2->{\textbf{18} against \textbf{0.6}: a gap of $H/u = 30\times$.}
    \end{column}
  \end{columns}
  \mlsyscredit{Trajectory Fine-Tuning, Theorem 16.1 (Ross and Bagnell 2010; Ross et al. 2011) and Napkin Math 16.3 (illustrative); values derived.}
  \note{[2.5 min] [CORE] ENGAGE: students predict both numbers before the curves appear. Define the terms first: $H$ is the number of decision steps; $\epsilon$ the chance the policy picks an action other than the reference's; cost per step is between 0 and 1; $u$ bounds the extra cost of one slip followed by recovery. Step 2: with error measured on the reference's states, the bound is $H^2\epsilon$: 900 times 0.02, 18 cost units of a possible 30. With the same rate measured on the learner's own states and $u = 1$: 0.6. Intuition for the square: the chance of a first divergence grows as $H\epsilon$, and a divergence can spoil the whole remaining horizon. Read in turns: up to 18 of 30 turns wasted, against fewer than one. To match the learner-state guarantee offline, the reference error would need to be 0.067 percent, 30 times smaller. The reference bound saturates at $H = 1/\epsilon = 50$ steps, beyond which it says nothing. LINK: Lecture 1's horizon compounding gave failure roughly $H\epsilon$; imitation on reference states adds a second factor of $H$. WARN: the linear bound needs $u$ small. One irreversible step (the remediation agent) makes $u$ grow with $H$ and restores the square; there the runtime's checks carry reliability. FLEX: [CORE].}
\end{frame}

\begin{frame}{DAgger labels the states the learner actually visits}
  \centering
  \mlsyssteps[width=0.94\textwidth,height=4.9cm,keepaspectratio]{dagger_step}{3}
  \mlsyscredit{Trajectory Fine-Tuning, figure 4(b) and \S16.6.2, redrawn; DAgger: Ross et al. 2011.}
  \note{[2.5 min] [CORE] Dataset Aggregation (DAgger). Step 1: run the current policy, mixed with the reference early on; collect the states it visits; ask a labeler (a larger model or a search procedure) what it should have done; execute that label in a sandbox and keep it only if the verifier passes; add the kept labels to the corpus; retrain; repeat. The labeler is now an online service in the collection loop, and a passing continuation shows that the label works while leaving open whether other actions also work. Step 2: fault injection supplies off-path states on purpose. The injected error is an observation with mask 0; the diagnosis and the corrective call are targets with mask 1. The selected targets inherit the fixture's causal meaning: a record that answers a lost acknowledgment with a blind retry teaches the duplicate effect that settlement before retry (Lecture 7) exists to prevent. Step 3: labeling learner states is what earns the linear bound. ENGAGE: what does fault injection cover that DAgger does not, and the reverse? (Expected: injection covers the failures designers anticipated; on-policy collection finds the ones the model's own behavior produces.) LINK: Lecture 17 drops the step-level label and scores fresh rollouts with a verifier. FLEX: [CORE].}
\end{frame}

\begin{frame}{In Super Tux Kart, DAgger stopped falling; cloning did not}
  \begin{columns}[c]
    \begin{column}{0.55\textwidth}
      \centering\safeimg[width=\textwidth,height=5.0cm,keepaspectratio]{stk.pdf}
    \end{column}
    \begin{column}{0.42\textwidth}
      \centering\safeimg[width=\textwidth,height=3.1cm,keepaspectratio]{supertuxkart_dagger.png}\\[0.08cm]
      {\footnotesize Star Track: a human steers; a linear policy imitates from the pixels.}
    \end{column}
  \end{columns}
  \mlsyscredit{Ross, Gordon, and Bagnell, AISTATS 2011, arXiv 1011.0686: Figure 1 (image) and Figure 2 (values read from the figure, redrawn) and \S5.1.}
  \note{[1.5 min] [OPTIONAL] Measured data from the DAgger paper. A human expert steered a kart around Super Tux Kart's Star Track, a track floating in space; the learner was a linear controller on image features, scored by falls off the track per lap. About 1,000 states were collected per iteration. Supervised cloning, which always trains on the expert's laps, did not improve as data grew: about 3 falls per lap after 1 iteration and about 4 after 20 (read from the figure). SMILe, an earlier mixing method, still fell about twice per lap after 20 iterations, as the paper's text states. DAgger, with the expert executed only in the first iteration, almost never fell after 5 iterations and never fell after 15. ENGAGE: why did more expert laps not help? (Expected: they repeat the same states; none shows how to recover once the kart drifts.) WARN: the values are read from a plot with 95 percent confidence bars; quote the shape, since the decimals are read by eye. FLEX: [OPTIONAL].}
\end{frame}

\begin{frame}{One static schema lets the weights memorize argument names}
  \begin{columns}[T]
    \begin{column}{0.50\textwidth}
      {\footnotesize\textbf{Proposed call against the delivered schema}}\par\vspace{0.05cm}
      \begin{logblock}[7.2]
        "arguments": \{\textcolor{RedLine}{\textbf{"file\_path"}}: "src/pool.py"\}\\[0.1cm]
        "current\_argument\_schema": \{\\
        \ind{2}"properties": \{\textcolor{GreenLine}{\textbf{"target\_uri"}}: \{"type": "string"\}\},\\
        \ind{2}"required": ["target\_uri"],\\
        \ind{2}"additionalProperties": false\}
      \end{logblock}
      \vspace{0.2cm}
      {\small The validator refuses it before dispatch.}
    \end{column}
    \begin{column}{0.47\textwidth}
      \small
      \textbf{Schema augmentation:} rename, reorder, add, or remove interface fields in \emph{every} part of the record together.\\[0.25cm]
      \uncover<2->{\textbf{Measure on renamed schemas:}\\
      SFA: calls that satisfy the delivered schema.\\
      PRE: calls that use a training-set key the schema lacks.}
    \end{column}
  \end{columns}
  \mlsyscredit{Trajectory Fine-Tuning, \S16.7, table 9, and \S16.11.2 (SFA, PRE); Gorilla (Patil et al. 2023).}
  \note{[2 min] [OPTIONAL] Coverage includes the interfaces the policy is shown. A corpus with one static tool signature lets the weights memorize argument names instead of reading the schema in context. The chapter's example: the runtime now requires target\_uri and forbids extra keys; a policy that memorized the older interface emits file\_path, and the validator refuses the call. Schema augmentation derives training records with declared interface changes: renamed tools or arguments, reordered fields, added unused tools, removed optional arguments. Every part of the record must change together: schema, selected calls, observations, and error messages. Renaming only the call while the old name survives in an earlier error message trains on a contradictory transcript. Step 2: measure it. Schema-following accuracy (SFA) is the share of calls that name a declared tool and satisfy its schema; parametric recall error (PRE) is the share that use a key from the training schemas that the delivered schema lacks. SFA rising on familiar schemas while PRE rises on renamed ones means the weights recall keys instead of reading the prompt. WARN: before blaming the weights, check that the serving template did not leave the old key in. Land the segment: where error is measured sets the horizon penalty, so long-horizon agents need labels on the states and interfaces they actually meet. FLEX: [OPTIONAL].}
\end{frame}

% =============================================================================
% SECTION: UPDATES
% =============================================================================
\mlsysnote{[0.5 min] [CORE] Fourth segment: ``fine-tuning'' names three independent choices: the objective, where the behavior comes from, and which weights change.}
\mlsysdivider{Which objective, which source, which weights?}{Trajectory Fine-Tuning, \S16.2.2 and \S16.8 to \S16.10}
\section{Updates}

\begin{frame}{When both repairs pass, a preference pair ranks them}
  \centering
  \safeimg[width=0.80\textwidth,height=3.1cm,keepaspectratio]{prefpair.pdf}\\[0.1cm]
  {\normalsize $\mathcal{L}_{\text{DPO}} = -\log \sigma\big(\beta\,[\,z_w - z_l\,]\big),\qquad z = \sum_{t\,\in\,\text{selected}} \log\dfrac{\pi_\theta(x_t \mid h_t)}{\pi_{\text{ref}}(x_t \mid h_t)}$}\\[0.15cm]
  \begin{columns}[c]
    \begin{column}{0.47\textwidth}
      \small Preference training by DPO. $\pi_{\text{ref}}$: frozen starting checkpoint. $\beta$: how strongly to resist drift from it.
    \end{column}
    \begin{column}{0.49\textwidth}
      \begin{logblock}[7.6]
        \# TRL 1.16.0, DPOConfig\\
        beta: float = field(default=0.1, ...\\
        precompute\_ref\_log\_probs: bool = ...\\
        \ind{2}"... the reference model does not need to be kept in memory."
      \end{logblock}
    \end{column}
  \end{columns}
  \mlsyscredit{Trajectory Fine-Tuning, \S16.8 and equation 7; Rafailov et al. 2023. Code: github.com/huggingface/trl, trl/trainer/dpo\_config.py, release 1.16.0 (Apache-2.0), condensed.}
  \note{[2.5 min] [CORE] Preference training. Two continuations of the connection-pool repair start from the same diagnosis and both fix the leak. One runs targeted tests and the full suite once; the other runs the full suite after every edit. Imitating both teaches the wasteful habit; a preference pair says which is better. Direct preference optimization (DPO) does it without a separate reward model. For each record, the score $z$ sums, over its selected policy tokens, the log ratio of the trained policy's probability to the frozen starting checkpoint's, $\pi_{\text{ref}}$, usually the SFT checkpoint. The loss pushes the preferred score $z_w$ above the dispreferred $z_l$; $\beta$ sets how hard the objective resists drifting from $\pi_{\text{ref}}$. Observations condition both records but carry no probability terms, and an identical shared prefix cancels exactly. Right: the real defaults: TRL sets $\beta = 0.1$ and can precompute the reference log-probabilities once, so the reference model need not sit in accelerator memory. WARN: e2 and e4 from the running case fail as a pair, because e4 saw the refusal and e2 did not. A pair needs identical available information. ENGAGE: what does $z_w - z_l > 0$ mean? (Expected: the policy moved toward the preferred record relative to the reference, nothing more.) FLEX: [CORE].}
\end{frame}

\begin{frame}{Is a 0.646 score the chance a repair succeeds?}
  \mlsysactivity{Vote}{2.5 min}
  \begin{columns}[T]
    \begin{column}{0.44\textwidth}
      \begin{mlsyscard}{mlsysaccent}
        \small Log-probability sums, policy / reference:\\
        preferred $-8$ / $-10$\\
        dispreferred $-11$ / $-10$\\[0.05cm]
        $z_w = 2$, $z_l = -1$, $\beta = 0.2$
      \end{mlsyscard}
      \vspace{0.15cm}
      \small\choice{A}{Yes, 64.6\%}\quad \choice{B}{No}
    \end{column}
    \begin{column}{0.52\textwidth}
      \small
      \renewcommand{\arraystretch}{1.2}
      \begin{tabular}{@{}lrr@{}}
        \toprule
        & $\beta = 0.2$ & $\beta = 0.4$ \\
        \midrule
        margin $\beta(z_w - z_l)$ & 0.6 & \uncover<2->{1.2} \\
        logistic score $\sigma$ & 0.646 & \uncover<2->{\textbf{0.769}} \\
        pair loss (from 0.693) & 0.437 & \uncover<2->{0.263} \\
        gradient coefficient & 0.0709 & \uncover<2->{0.0926} \\
        \bottomrule
      \end{tabular}\\[0.2cm]
      \uncover<2->{\textbf{B.} The same log-probabilities give a higher score when $\beta$ doubles.}
    \end{column}
  \end{columns}
  \mlsyscredit{Trajectory Fine-Tuning, Napkin Math 16.4 (hypothetical log-probabilities: illustrative) and \S16.13; values derived in images/calcs.py.}
  \note{[2.5 min] [CORE] ENGAGE: vote A or B, one reason. Walk the left column: the scores are $-8 - (-10) = 2$ and $-11 - (-10) = -1$; at $\beta = 0.2$ the margin is 0.6, the logistic score 0.646, and the pair loss 0.437, down from $\ln 2 = 0.693$ at the start. The saturation factor $\sigma(-\text{margin})$ is 0.354, so the full gradient coefficient is $\beta$ times that, 0.0709. Step 2: double $\beta$ with the same log-probabilities: the score rises to 0.769 and the loss falls to 0.263, while the gradient coefficient rises to 0.0926. A quantity that moves with a hyperparameter measures fit to this one comparison and says nothing about the repair's success rate. And the margin can widen by pushing the dispreferred record down while the preferred one stays put. The answer is B; acceptance is measured in closed-loop trials. WARN: the Bradley--Terry model behind DPO identifies only differences of scores, so a single record's score has no meaning outside its pair. FLEX: [CORE].}
\end{frame}

\begin{frame}{Distilling a checklist pays only if it changes slower than the bill}
  \begin{columns}[c]
    \begin{column}{0.50\textwidth}
      \centering\safeimg[width=\textwidth,height=4.9cm,keepaspectratio]{distill.pdf}
    \end{column}
    \begin{column}{0.47\textwidth}
      \small
      \textbf{Context distillation:} train the model to act without a 3,000-token checklist it used to read every turn.\\[0.15cm]
      \textbf{Teacher distillation:} imitate a stronger model's \emph{verified} trajectories.\\[0.25cm]
      \uncover<2->{Uncached: \textbf{\$33,000} a month net. With the checklist cached: \textbf{\$10,200} a month lost.}
    \end{column}
  \end{columns}
  \mlsyscredit{Trajectory Fine-Tuning, \S16.9 and Exercise 8 (support agent; hypothetical prices and volumes: illustrative); savings derived in images/calcs.py.}
  \note{[2.5 min] [CORE] Trajectory distillation: a costly teacher configuration produces trajectories that pass the same admission as any demonstration, and a cheaper student imitates them. Two forms. Context distillation: the teacher is the same model with a long procedural prompt; the loader strips the guidance and trains the student to emit the same admitted actions from the short prompt. Teacher distillation: the teacher is a stronger model or a search harness reached over an API, which returns sampled tokens without probability distributions, so distillation is filtered imitation of verified trajectories; soft targets would need the distributions. The worked case, from the chapter's exercise for the support agent: a 3,000-token checklist read on every turn, 4 turns per session, 2 million sessions a month, 24 billion input tokens at \$2 per million: \$48,000 a month. Redistilling after each monthly checklist revision costs \$15,000, so the net is \$33,000. Step 2: if the stable prefix is cached at one tenth of the price, the checklist costs only \$4,800 a month, and distilling loses \$10,200. ENGAGE: at how many revisions a month does the uncached saving vanish? (Expected: 3.2.) WARN: latency savings must be measured on closed-loop runs; one extra recovery turn erases the prefill saving. LINK: Lecture 19 prices adaptation payback in general. FLEX: [CORE].}
\end{frame}

\begin{frame}{LoRA trains two thin matrices beside each frozen one}
  \begin{columns}[T]
    \begin{column}{0.46\textwidth}
      \centering
      {\normalsize $W = W_0 + \dfrac{\alpha}{r}\,B A$}\\[0.15cm]
      {\small $B$: $d_{\text{out}}\times r$,\quad $A$: $r\times d_{\text{in}}$,\quad $r \ll d$}\\[0.1cm]
      {\small Parameters per matrix: $r\,(d_{\text{in}} + d_{\text{out}})$}\\[0.3cm]
      {\small $32 \times 2 \times 16 \times (4{,}096 + 4{,}096)$\\$= \mathbf{8{,}388{,}608}$}
    \end{column}
    \begin{column}{0.51\textwidth}
      {\footnotesize\textbf{Our run: PEFT on the Llama 3 8B architecture}}\par\vspace{0.05cm}
      \begin{mlsyscard}{midgray}
      \begin{logblock}[9.5]
        LoraConfig(r=16, lora\_alpha=32,\\
        \ind{2}target\_modules=['q\_proj', 'o\_proj'])\\[0.15cm]
        trainable params: \colorbox{mlsysaccent!12}{\textbf{8,388,608}} ||\\
        all params: 8,038,649,856 ||\\
        trainable\%: \textbf{0.1044}
      \end{logblock}
      \end{mlsyscard}
      \vspace{0.2cm}
      {\small The chapter's hypothetical matches a real architecture.}
    \end{column}
  \end{columns}
  \mlsyscredit{Trajectory Fine-Tuning, \S16.10, Napkin Math 16.5; Hu et al. 2021. Run: peft 0.21.2 on the Meta-Llama-3-8B config.json, no weights (images/tool\_runs.py).}
  \note{[2 min] [CORE] Parameter-efficient fine-tuning freezes the base and trains a small adapter. Low-rank adaptation (LoRA) writes the update to a chosen matrix $W_0$ as the product of two thin matrices, $B$ with $r$ columns and $A$ with $r$ rows, scaled by $\alpha/r$. $B$ starts at zero, so training starts from exactly the base model. Each adapted matrix adds $r(d_{\text{in}} + d_{\text{out}})$ parameters. The chapter's case: rank 16 on two 4,096 by 4,096 matrices in each of 32 layers, 8,388,608 parameters, 0.10 percent of 8 billion. Right: we ran the real library on the real Llama 3 8B architecture (the configuration file only, on PyTorch's meta device, so no weights were downloaded). In Llama 3 8B the query and output projections are both 4,096 by 4,096, and PEFT reports exactly 8,388,608 trainable parameters, 0.1044 percent of 8.04 billion. WARN: not every projection is square; Llama 3's key and value projections are 4,096 by 1,024, so adapting them adds fewer parameters. Sum $r(d_{\text{in}} + d_{\text{out}})$ over the actual targets. ENGAGE: does training 0.1 percent of the parameters cut training memory to about 0.1 percent? Hold the vote for the next slide. FLEX: [CORE].}
\end{frame}

\begin{frame}{How much memory does training 0.10\% keep?}
  \mlsysactivity{Predict}{2 min}
  \begin{columns}[c]
    \begin{column}{0.50\textwidth}
      \centering\uncover<2->{\safeimg[width=\textwidth,height=4.9cm,keepaspectratio]{loramem.pdf}}
    \end{column}
    \begin{column}{0.47\textwidth}
      \small
      \choice{A}{about 0.1\%}\quad \choice{B}{about 1\%}\quad \answerat{2}{C}{about 13\%}\\[0.3cm]
      \uncover<2->{128 GB $\rightarrow$ 16.13 GB: \textbf{7.9$\times$}, never more than 16 B / 2 B $=$ 8$\times$.\\[0.2cm]}
      \uncover<3->{\textbf{Measured, GPT-3 175B:} training memory 1.2 TB $\rightarrow$ 350 GB (3.4$\times$); 100 adapted copies 354 GB against 35 TB.}
    \end{column}
  \end{columns}
  \mlsyscredit{Trajectory Fine-Tuning, Napkin Math 16.5 (hypothetical 8B model, mixed-precision Adam: illustrative). Measured: Hu et al., LoRA, arXiv 2106.09685, \S4.2.}
  \note{[2 min] [CORE] ENGAGE: vote first. Step 2: C. Under mixed-precision Adam each trained parameter carries 16 bytes (2 weight, 2 gradient, 4 master copy, 4 and 4 for the two moments), and each frozen parameter only its 2-byte weight. Full fine-tuning of 8 billion parameters: 128 GB of parameter state. The LoRA run: 16 GB of frozen base plus 134.22 MB of adapter state, 16.13 GB, 12.6 percent of full, a 7.9 times reduction. The frozen base is 99.2 percent of what remains, so no choice of rank can push the saving past 16 over 2, 8 times; doubling the rank adds only another 134 MB. The lever is the frozen base's precision, which quantized-base methods such as QLoRA exploit. Activations come on top for both. Step 3: measured in the LoRA paper on GPT-3 175B: training memory fell from 1.2 TB to 350 GB, 3.4 times, and 100 task adaptations stored as adapters take about 354 GB instead of 35 TB. WARN: the memory fraction is far larger than the parameter fraction. FLEX: [CORE].}
\end{frame}

\begin{frame}{An adapter is a set of weights, pinned to its base and granted nothing}
  \begin{columns}[c]
    \begin{column}{0.47\textwidth}
      {\footnotesize\textbf{adapter\_config.json, as PEFT wrote it (excerpt)}}\par\vspace{0.05cm}
      \begin{mlsyscard}{midgray}
      \begin{logblock}[9.5]
        "base\_model\_name\_or\_path":\\
        \ind{2}"meta-llama/Meta-Llama-3-8B",\\
        \colorbox{RedLine!12}{\textcolor{RedLine}{\textbf{"revision": null,}}}\\
        "r": 16,\\
        "target\_modules": ["q\_proj", "o\_proj"]
      \end{logblock}
      \end{mlsyscard}
    \end{column}
    \begin{column}{0.50\textwidth}
      \centering
      \begin{tikzpicture}[font=\small, every node/.style={draw, rounded corners=3pt, line width=1pt, align=center, minimum height=0.95cm}]
        \node[fill=BlueLine!10, draw=BlueLine, minimum width=2.4cm] (ad) at (0,0) {\textbf{Adapter} $BA$};
        \node[fill=RedLine!8, draw=RedLine, minimum width=2.4cm] (base) at (3.85,0) {\textbf{Base} $W_0$\\exact revision};
        \node[fill=mlsysaccent!8, draw=mlsysaccent, minimum width=6.1cm] (rt) at (1.93,-2.0) {\textbf{Runtime}: grants, limits, caches};
        \draw[-{Stealth[length=6pt]}, line width=1.2pt, RedLine] (ad) -- node[draw=none, above=1pt, font=\footnotesize, minimum height=0pt, inner sep=1pt] {pinned to} (base);
        \draw[-{Stealth[length=6pt]}, line width=1.2pt, mlsysaccent] (ad) -- node[draw=none, left, font=\footnotesize, minimum height=0pt] {granted nothing} (ad |- rt.north);
      \end{tikzpicture}\\[0.35cm]
      {\small \textbf{Pin the base}; share \textbf{no prefix cache} across adapters; grant capabilities \textbf{per role}.}
    \end{column}
  \end{columns}
  \mlsyscredit{Our run: peft 0.21.2 LoraConfig.save\_pretrained (images/tool\_runs.py, data/adapter\_config.json), 5 of 41 keys shown. Trajectory Fine-Tuning, \S16.10.2.}
  \note{[1.5 min] [OPTIONAL] The real release record PEFT writes beside every adapter, five of its keys. The diagram: the delta $BA$ belongs to one exact base $W_0$, and the runtime, not the adapter, holds grants and limits. The record names the base model, but the revision field is null by default: nothing pins which version of the base the delta was trained against. On a silently updated base the adapter still loads, and its outputs are wrong. The chapter's release manifest pins the base version and every training input. Two serving rules follow. Cached prefix state computed under one adapter cannot serve another adapter that changed those layers (Lecture 5 owns prefix-cache reuse). And serving several agent roles as adapters over one shared base saves memory, but the roles share one process, its failure domain, and its limits; an adapter holds no permissions of its own, so each role still needs its own capability grants under least privilege (Lecture 8). ENGAGE: a single role in production: keep the adapter separate or fold it in? (Expected: folding $W_0 + \frac{\alpha}{r}BA$ into one dense matrix removes adapter overhead per token when only one role is served.) Land the segment: objective, source, and artifact are independent; each changes the cost and the evidence the candidate needs. FLEX: [OPTIONAL].}
\end{frame}

% =============================================================================
% SECTION: RELEASE
% =============================================================================
\mlsysnote{[0.5 min] [CORE] Last segment: the trained candidate exists. What evidence promotes it?}
\mlsysdivider{When does a fine-tuned policy earn a release?}{Trajectory Fine-Tuning, \S16.11 to \S16.13}
\section{Release}

\begin{frame}{Should the lower-loss checkpoint ship?}
  \mlsysactivity{Predict}{2 min}
  \begin{columns}[c]
    \begin{column}{0.48\textwidth}
      \small
      Same curated trajectories, two runs:\\[0.15cm]
      \textbf{Run 1:} loss on all tokens. \textcolor{GreenLine}{\textbf{Lower dashboard loss.}}\\
      \textbf{Run 2:} loss on selected targets.\\[0.3cm]
      \choice{A}{promote run 1}\quad \choice{B}{promote run 2}\\[0.1cm]
      \answerat{2}{C}{neither, yet}
    \end{column}
    \begin{column}{0.49\textwidth}
      \small
      \uncover<2->{Run 1 averages over 16,383 terms, 77.1\% observations; run 2 over 800 targets. Different units.\\[0.2cm]
      Recompute held-out loss under one mask, split, template, and normalization; rank offline.\\[0.2cm]
      \textbf{Promote on closed-loop acceptance} under the serving runtime.}
    \end{column}
  \end{columns}
  \mlsyscredit{Trajectory Fine-Tuning, Example 16.1 and \S16.11.1; counts from Napkin Math 16.1 (illustrative).}
  \note{[2 min] [CORE] ENGAGE: one minute; vote. Step 2: C. The two losses average over different token sets. In Napkin Math 16.1, the all-token loss averages 16,383 prediction terms, most of them observations whose repetitive parts (echoed schemas, boilerplate, repeated paths) are easy to predict; the masked loss averages 800 action tokens. Comparing them rewards run 1 for predicting compiler output. Even a held-out loss recomputed under the same mask, split, tokenizer, template, and normalization only ranks candidates offline: it is measured on teacher-forced prefixes, and the policy's own rollouts can still call invalid tools, misread an unexpected exception, or exhaust the budget. Promotion rests on budgeted task acceptance under the task acceptance contract of Lecture 13, as paired trials against the incumbent. WARN: if two checkpoints use different tokenizers, their per-token losses do not even share a unit. FLEX: [CORE].}
\end{frame}

\begin{frame}{Higher acceptance: should this candidate ship?}
  \mlsysactivity{Diagnose}{2.5 min}
  \begin{columns}[T]
    \begin{column}{0.52\textwidth}
      \footnotesize
      \renewcommand{\arraystretch}{1.2}
      \begin{tabular}{@{}lrr@{}}
        \toprule
        500 trials each & Baseline & Candidate \\
        \midrule
        accepted & 60\% & 66\% \\
        mean attempt cost & \$1.20 & \$1.50 \\
        $C_{\text{eff}}$ & \uncover<2->{\$2.00} & \uncover<2->{\textbf{\$2.27}} \\
        SFA, familiar schemas & 0.92 & 0.95 \\
        SFA, renamed schemas & 0.90 & \textcolor{RedLine}{\textbf{0.81}} \\
        PRE, renamed schemas & 0.02 & \textcolor{RedLine}{\textbf{0.12}} \\
        \bottomrule
      \end{tabular}
    \end{column}
    \begin{column}{0.45\textwidth}
      \small
      \uncover<2->{Cost per accepted task rises 14\%; break-even attempt cost \$1.32.\\[0.15cm]
      Difference: 1.97 standard errors, if trials were independent.\\[0.15cm]}
      \uncover<3->{\textbf{Hold it.} PRE rose sixfold, so the weights recall old keys. Test schema augmentation first; rerun paired, with retention.}
    \end{column}
  \end{columns}
  \mlsyscredit{Trajectory Fine-Tuning, Exercise 9 (hypothetical: illustrative); $C_{\text{eff}}$, break-even, and standard error derived in images/calcs.py.}
  \note{[2.5 min] [OPTIONAL] ENGAGE: pairs, 90 seconds: name one reason to ship and two reasons not to. Step 2: cost per accepted task, the attempt cost divided by the acceptance rate (Lecture 1): baseline \$1.20 / 0.60 = \$2.00; candidate \$1.50 / 0.66 = \$2.27, 14 percent more. The candidate would match the baseline only at a mean attempt cost of \$1.32. Treating the arms as independent, the standard error of the difference is 0.0305, and the 6-point gain is 1.97 standard errors, borderline. If both arms ran on the same 500 tasks, the paired comparison of Lecture 13 needs the per-task outcomes, and the count of tasks where only one policy succeeded decides. Step 3: the decisive row is the renamed-schema set. SFA fell from 0.90 to 0.81 and PRE rose from 0.02 to 0.12: the weights recall the training schema's keys. The decision is hold; the first training change to test is schema augmentation, and the gate needs repeated paired trials, the version-conflict state, and a retention suite. WARN: the all-token training loss that fell is no evidence at all. FLEX: [OPTIONAL].}
\end{frame}

\begin{frame}{Vote: true or false?}
  \mlsysactivity{Vote}{2 min}
  \small
  \renewcommand{\arraystretch}{1.1}
  \begin{tabular}{@{}>{\raggedright\arraybackslash}m{0.47\textwidth}>{\raggedright\arraybackslash}m{0.47\textwidth}@{}}
    \toprule
    \textbf{Claim} & \uncover<2->{\textbf{Verdict}} \\
    \midrule
    Every response in an accepted trajectory is a good target. & \uncover<2->{\textbf{\textcolor{RedLine}{False.}} The refused call is context: 160 of 960 positions.} \\
    Packing changes only speed. & \uncover<2->{\textbf{\textcolor{RedLine}{False.}} Masks and normalization change the lesson.} \\
    Recovery examples make the agent robust to the states it will meet. & \uncover<2->{\textbf{\textcolor{RedLine}{False.}} 18 against 0.6: label the learner's states.} \\
    Promote the checkpoint with the lowest loss. & \uncover<2->{\textbf{\textcolor{RedLine}{False.}} Promote on closed-loop acceptance and $C_{\text{eff}}$.} \\
    \bottomrule
  \end{tabular}
  \mlsyscredit{Trajectory Fine-Tuning, \S16.13 (fallacies and pitfalls) with Napkin Math 16.1 and 16.3 (illustrative).}
  \note{[2 min] [CORE] ENGAGE: students write T or F for each, then compare with a neighbor. Step 2: all four are the chapter's fallacies and pitfalls, and all are false. One: the refused proposal is the one action the policy must never learn to emit; an unselective mask spends 160 of 960 response positions on it. Two: an ordinary causal mask across a packed buffer leaks context, and token against example averaging changes weights 32 times in the worked batch. Three: recovery records and perturbed schemas broaden coverage, but the bound on reference states is 18 against 0.6 on learner states, and the linear bound holds only while one slip stays cheap. Four: lower loss may only mean better prediction of tool output. The two fallacies not on the table, a preference score as a success probability and parameter fraction as memory fraction, were the two votes of the last segment. FLEX: [CORE].}
\end{frame}

\begin{frame}{In the real runs, only part of the model's 12 to 31\% is a target}
  \begin{columns}[c]
    \begin{column}{0.56\textwidth}
      \centering\safeimg[width=\textwidth,height=4.6cm,keepaspectratio]{roles_col_step2.pdf}
    \end{column}
    \begin{column}{0.41\textwidth}
      \small
      \textcolor{GreenLine}{\textbf{Green}}: context, mask 0, still read.\\[0.1cm]
      \textbf{Dark}: only admitted spans, with their stop tokens.\\[0.1cm]
      Batches: declared weighting, isolated examples.\\[0.1cm]
      Coverage: label the states the policy reaches.\\[0.1cm]
      Release: closed-loop acceptance, runtime unchanged.
    \end{column}
  \end{columns}
  \mlsyscredit{SWE-bench experiments trajectories via MLSysIM (as in the opening); Trajectory Fine-Tuning, \S16.12 and \S16.14.}
  \note{[2.5 min] [CORE] Return to the opening bars. The model wrote 12 to 31 percent of these real transcripts. The lecture's answer to the driving question, why can a useful recovery record teach the wrong lesson: because the loss teaches whatever it touches. Green, the observations, stays in the input with mask 0 and is still read. Of the dark region, only the spans curation admitted are targets, each with its authentic stop token; the refused call is context. The batch declares its weighting and keeps packed examples isolated. The corpus covers only recorded states, so the policy needs labels on the states it reaches. The chapter's case study settles trained\_policy\_v1 this way: imitation objective, because the record supplies a correction and no admissible comparison; the admitted recovery corpus as the source; a full-parameter update as the reference any later adapter must match, at 128 GB against 16.13 GB. And whatever the training did, the runtime keeps validating schemas, enforcing least privilege, and bounding time; the final report is still a claim the harness checks. ENGAGE: which green region in these runs would you be most tempted to train on, and why not? (Expected: echoed file contents; the model would learn to invent them.) FLEX: [CORE].}
\end{frame}

\begin{frame}{Takeaways}
  \thumblineL{thumb_homes.pdf}{Weights hold procedure; context holds contracts; the runtime enforces.}\\[0.06cm]
  \thumblineL{thumb_mask.pdf}{Train only selected policy spans; unmasked, targets dilute 20.5$\times$.}\\[0.06cm]
  \thumblineL{thumb_pack.pdf}{Declare and test normalization and packing (32$\times$ weighting).}\\[0.06cm]
  \thumblineL{thumb_bound.pdf}{Measure error where the learner goes: 18 against 0.6.}\\[0.06cm]
  \thumblineL{thumb_pref.pdf}{A preference is one comparison; LoRA cuts state only 7.9$\times$.}\\[0.06cm]
  \thumblineL{thumb_gate.pdf}{Gate release on closed-loop acceptance and $C_{\text{eff}}$.}
  \note{[1 min] [CORE] One line per goal; the full statements: weights hold procedure, schemas and permissions stay in context, and enforcement stays in the runtime. Only selected policy spans are targets; unmasked, each target is diluted 20.5 times and the policy learns to invent tool results. Normalization and packing can change the lesson through 32 times weighting and cross-example leaks. Where error is measured sets the horizon penalty: 18 against 0.6 at 2 percent over 30 steps. A preference score measures one comparison; LoRA trains a small share of parameters yet cuts state only 7.9 times in bytes. Ask students which slide taught each line. If short on time, read the second and the last.}
\end{frame}

\begin{frame}{Open questions}
  \normalsize
  \begin{itemize}\setlength{\itemsep}{0.3cm}
    \item Rationale tokens: train the model to write its own reasoning, or keep reasoning as context? What would you measure to decide?
    \item DAgger needs a labeler on call for every visited state. When is that labeler more expensive than the failures it prevents?
    \item Schemas change weekly in many deployments. Should agents be fine-tuned at all on tool use, or only on procedure?
  \end{itemize}
  \note{[1 min] [OPTIONAL] No single answer. Pick one for discussion or assign one as a short written response. The first comes from the chapter's three rationale treatments (\S16.4.4): selected, kept as context, or removed; each asks a different serving question. The second asks students to price on-policy labeling with cost per accepted task, as Exercise 5 does. The third pushes on the three homes: the procedure of reading a schema can be trained; the schema itself cannot.}
\end{frame}

\begin{frame}{Read and explore}
  \small
  \begin{itemize}\setlength{\itemsep}{0.04cm}
    \item \textbf{Read:} Trajectory Fine-Tuning, \S16.2 to \S16.12; Checkpoints 16.1 and 16.2; Exercises 2, 3, and 7.
    \item Ross, Gordon, and Bagnell, ``A Reduction of Imitation Learning \ldots{} (DAgger),'' AISTATS 2011.
    \item Rafailov et al., ``Direct Preference Optimization,'' NeurIPS 2023.
    \item Hu et al., ``LoRA: Low-Rank Adaptation of Large Language Models,'' 2021.
    \item Krell et al., ``Efficient Sequence Packing Without Cross-Contamination,'' 2021.
    \item Hugging Face, ``Fixing Gradient Accumulation,'' 2024.
  \end{itemize}
  \note{[0.5 min] [CORE] The chapter first, with its two checkpoints; Exercises 2, 3, and 7 practice dilution, normalization, and adapter memory by hand. Each outside item backs one segment: the regret bounds and the Super Tux Kart data (Ross), the preference loss and which assumptions a span surrogate relaxes (Rafailov), low-rank updates and the measured memory savings (Hu), packing as an equivalence to test (Krell), and the real accumulation incident (Hugging Face). FLEX: add the training library your course uses.}
\end{frame}

\begin{frame}[c]
  \centering
  {\small\textcolor{lightgray}{NEXT LECTURE \textbar{} REINFORCEMENT LEARNING FROM VERIFIABLE REWARDS}}\\[0.4cm]
  \begin{columns}[c]
    \begin{column}{0.30\textwidth}
      \centering\safeimg[width=\textwidth,height=1.6cm,keepaspectratio]{thumb_next.pdf}
    \end{column}
    \begin{column}{0.66\textwidth}
      {\LARGE\bfseries What does an agent learn when the check becomes the reward?}
    \end{column}
  \end{columns}
  \note{[0.5 min] [CORE] Every update today needed a demonstrated or compared response, so the version conflict trained\_policy\_v1 never saw could be taught only by labeling or injecting it. Lecture 17 starts from that checkpoint and lets the acceptance check score fresh attempts on the states the policy reaches itself. The check stops being only a release gate and becomes the objective, which raises a new question: what the optimizer does with every gap between the check and the task.}
\end{frame}

\end{document}
